This thesis investigates whether higher order pattern unification can be
integrated with a type checker that represents name binding using scope graphs.
The work is carried out in Statix, a constraint language in the Spoofax
language workbench. Statix uses scope graphs for declarative name resolution
and provides first order unification as part of its constraint solver.
However, inference in dependently typed languages can require higher order
unification when the inferred solution depends on variables introduced by
surrounding binders.

We implement a small dependently typed language in Statix and extend its type
checker with higher order pattern unification. Scope graph information is used
to determine which bound variables may occur in a solution and to construct
well scoped higher order solutions. The implementation supports the pattern
cases considered in this thesis, but does not support postponement of
constraints that may become solvable later.

We also implement the core of the same language in Elpi using Higher Order
Abstract Syntax. The two implementations are used to compare explicit
scope based binding with binding provided by a higher order meta language.
The comparison shows that Elpi provides direct support for lexical binding,
substitution, and higher order unification, while scope graphs provide a
direct representation of non lexical name resolution such as modules,
qualified references, and forward references.